\documentclass{article}
\usepackage{multicol}
\usepackage{/globals/math}
\begin{document}

\section{Binomische Formeln}
Werden zwei Binome (Polynome bestehend aus zwei Gliedern) gleicher Glieder mit einander multipliziert, so können diese mit den binomischen Formeln umgeformt werden.

Diese sind keine neuen Regeln an sich; sie sind einfache generalisierungen des Ausmultiplizieren.
\begin{mathdescribe}
 \mathitem{1. binomische Formel} (a+b)^2 = a^2+2ab+b^2\\ 
 \mathitem{2. binomische Formel} (a-b)^2 = a^2-2ab+b^2 \\
 \mathitem{3. binomische Formel} (a+b)(a-b) = a^2-b^2
\end{mathdescribe}

\subsection{Höhere Potenzen}
Die ersten beiden binomischen Formeln können auch generalisiert werden, auf Binome mit anderen Exponenten außer 2. Dies kann rein algebraisch passieren
\[
 (a+b)^n = \sum_{k=0}^n \binom{n}{k} a^{n-k}b^k
\]
Ansonsten kann hier auch das Pascalsche Dreieck helfen. Hier wird ein Dreieck gebildet, indem als erste Reihe $1 \qquad 1$ aufgeschrieben wird und für die nächsten Reihen zwei nebeneinanderstehende Zahlen zusammenaddiert werden und an den Rändern neue einsen aufgeschrieben werden.

Wird sich die $n$te Reihe angeguckt und die $k$te Zahl in dieser gefunden, so ist diese $\binom{n}{k}$.
\end{document}